\documentclass[10pt,a4paper]{amsart}
\usepackage[margin=19mm]{geometry}
\usepackage{amsmath,amssymb,mathtools}
\usepackage[scr=rsfs,cal=euler]{mathalfa}
\usepackage{xfrac}
\usepackage[unicode,hidelinks,bookmarks=false]{hyperref}
\newcommand{\ie}{i.e.\ }
\newcommand{\cf}{cf.\ }
\newcommand{\resp}{respectively\ }
\newcommand{\Subsection}{\subsection}
\providecommand{\nobreakdash}{\nobreak\hbox{-}}
\setlength{\emergencystretch}{2em}
\multlinegap0pt
\usepackage{a4wide}
\usepackage{lmodern}
\usepackage{mathtools}
\newtheorem{theorem}{Theorem}[section]
\newtheorem{proposition}[theorem]{Proposition}
\newtheorem{lemma}[theorem]{Lemma}
\newtheorem*{theoremA}{Theorem A}
\newtheorem*{corollaryB}{Corollary B}
\newcommand{\R}{\mathbb{R}}
\newcommand{\N}{\mathbb{N}}
\newcommand{\abs}[1]{\lvert #1\rvert}
\newcommand{\norm}[1]{\lVert #1\rVert}
\newcommand{\Fin}{\mathcal F}
\renewcommand{\bibliofont}{\scriptsize}
\begin{document}
\title[Pointwise function spaces over compacta are not weak Banach ]{Pointwise function spaces over compacta are not weak Banach spaces}
\author{Tomasz Kania; Jerzy Kąkol}
\date{}
\maketitle
\noindent\emph{Source and licence.} Pointwise function spaces over compacta are not weak Banach spaces. Version arXiv:2608.17549v1 (CC BY 4.0). This material is distributed under the Creative Commons Attribution 4.0 International licence, http://creativecommons.org/licenses/by/4.0/, as recorded in the arXiv OAI licence metadata for this exact version and shown on the abstract page https://arxiv.org/abs/2608.17549v1. Full rights notice: licenses/arxiv-2608.17549-CC-BY-4.0.txt.\par\smallskip
\noindent\textit{Editorial scope.} Counted unit: the original \texttt{lemma} environment \texttt{lem:fixed-support} at source lines 315--325 together with its complete proof at lines 327--351; the delivered document keeps source lines 162--352 so that every referenced label and every citation resolves inside it. Native editable source is the paper's own concatenated TeX; the AMS wrapper is editorial formatting only. No publisher class, upstream script or unrelated artwork is redistributed.
\begin{lemma}\label{lem:arc}
Let $E$ be an infinite-dimensional real Banach space.  There is a
norm-continuous map $[0,1]\to B_{E^*}$, $t\mapsto\mu_t$,
whose image $P$ is homeomorphic to $[0,1]$, does not contain $0$, and is
linearly independent.  Consequently, if $V\subseteq E^*$ is
linear and $\dim V=r<\infty$, then $\abs{V\cap P}\leqslant r$.  The norm
and weak-star topologies coincide on $P$.
\end{lemma}

\begin{proof}
The dual $E^*$ is infinite-dimensional.  By Mazur's basic-sequence
selection theorem, it contains a normalised basic sequence
$(\nu_j)_{j\geqslant0}$; see
\cite[Section~1.5]{AlbiacKalton2016}.  We use the convention $t^0=1$
for every $t\in[0,1]$, including $t=0$.  For $t\in[0,1]$, let
$\mu_t=\sum_{j=0}^{\infty}2^{-j-1}t^j\nu_j$.
The series converges absolutely in norm and $\norm{\mu_t}\leqslant1$.
Moreover, the terms corresponding to $j=0$ cancel in
$\mu_t-\mu_s$, and the mean value theorem gives
$\abs{t^j-s^j}\leqslant j\abs{t-s}$ for $j\geqslant1$ and
$s,t\in[0,1]$.  Therefore
$\norm{\mu_t-\mu_s}\leqslant
\sum_{j=1}^{\infty}2^{-j-1}\abs{t^j-s^j}\leqslant
\abs{t-s}\sum_{j=1}^{\infty}j2^{-j-1}=\abs{t-s}$.  Thus the map is
norm-continuous.

Let us suppose that $t_1,\ldots,t_q$ are distinct and
$a_1\mu_{t_1}+\cdots+a_q\mu_{t_q}=0$.
Uniqueness of the coefficients with respect to the basic sequence gives
$a_1t_1^j+\cdots+a_qt_q^j=0$ for every $j\geqslant0$.  For
$j=0,\ldots,q-1$, the coefficient matrix is the transpose of a
Vandermonde matrix.  Since the $t_i$ are distinct, it is nonsingular by
the standard determinant formula
\cite[Section~VM, Theorem~DVM]{Beezer2004}, so
$a_1=\cdots=a_q=0$.  Hence $P=\{\mu_t:0\leqslant t\leqslant1\}$ is
linearly independent.  In particular the map is injective, and
the non-zero coefficient of $\nu_0$ shows that $0\notin P$.  A continuous
injection from the compact interval into the normed space $E^*$ is a
homeomorphic embedding, so $P$ is a norm-compact arc and is perfect.

Any subset of an $r$-dimensional space having more than $r$ elements is
linearly dependent, which proves the stated estimate.  Finally,
the identity from $P$ with its norm topology to $P$ with its weak-star
topology is a continuous bijection from a compact space to a Hausdorff
space.  It is therefore a homeomorphism.
\end{proof}

The next lemma is the countable $k$-uniform instance of the
$\Delta$-system theorem of Erd\H{o}s and Rado
\cite[p.~86]{ErdosRado1960}.  We include its short proof in the form
used below.

\begin{lemma}\label{lem:delta-system}
Let $k\in\N$ and let $(F_n)$ be a sequence of pairwise distinct
$k$-element sets.  There are strictly increasing indices
$n_1<n_2<\cdots$, a finite set $A$, and pairwise disjoint non-empty
finite sets $A_j$, disjoint from $A$, such that
$F_{n_j}=A\cup A_j$ for every $j\in\N$.  In particular, $\abs A<k$.
The set $A$ is called the root of this $\Delta$-system, and the sets
$A_j=F_{n_j}\setminus A$ are called its petals.
\end{lemma}

\begin{proof}
We argue by induction on $k$.  The case $k=1$ is immediate, with
$A=\varnothing$.  Let us suppose that $k>1$.  If some point $x$ belongs to
infinitely many $F_n$, let us retain those sets and apply the induction
hypothesis to the pairwise distinct sets $F_n\setminus\{x\}$; then
let us adjoin $x$ to the root.  If no point belongs to infinitely many $F_n$,
let us choose a pairwise disjoint subsequence recursively.  Indeed, a finite
union of previously chosen sets meets only finitely many $F_n$.  In the
latter case, let us take $A=\varnothing$.
\end{proof}

\section{Proof of the main theorem}
\label{sec:main-proof}

Let us suppose, towards a contradiction, that a map as in Theorem~A exists.
Translation in $E_w$ allows us to assume that
$h\colon C_p(K)\to E_w$ satisfies $h(0)=0$,
while preserving continuity of $h$ and continuity of $h^{-1}$ at the
origin.  If $K=\varnothing$, surjectivity is already impossible, so let
us assume that $K$ is non-empty.  Let us fix the arc
$P\subseteq B_{E^*}$ supplied
by Lemma~\ref{lem:arc}.
In what follows, $P$ carries its norm topology, which coincides with
its weak-star topology by Lemma~\ref{lem:arc}.

For $k\in\N$, let
$\Fin_k(K)=\{F\subseteq K:1\leqslant\abs F\leqslant k\}$ carry the
Vietoris topology, and let $\Fin_0(K)=\varnothing$.  The map from $K^k$
onto $\Fin_k(K)$ which sends $(x_1,\ldots,x_k)$ to
$\{x_1,\ldots,x_k\}$ is continuous: the inverse image of a basic
Vietoris set $\langle U_1,\ldots,U_s\rangle$ consists of those tuples
whose coordinates lie in $\bigcup_{i=1}^s U_i$ and meet every $U_i$,
and is therefore open.  Hence $\Fin_k(K)$ is compact.  Since $K$ is
Hausdorff, its Vietoris hyperspace, and therefore $\Fin_k(K)$, is
Hausdorff.  For a
finite $F\subseteq K$ and $m\in\N$, let
$U(F,m)=\{f\in C(K):\abs{f(x)}<1/m\text{ for }x\in F\}$; in particular,
$U(\varnothing,m)=C_p(K)$.  For $\mu\in P$, let
$H(\mu)=\{y\in E:\abs{\mu(y)}\leqslant1\}$.

The following relation is adapted from the relation in the proof of
\cite[Theorem~4.1]{KrupskiMarciszewski2017}.  Let
$\pi_2\colon\Fin_k(K)\times P\to P$ denote the projection onto the
second coordinate, so that $\pi_2(F,\mu)=\mu$.  Let
\[
 \begin{split}
 Z_{k,m}&=\{(F,\mu)\in\Fin_k(K)\times P:
                    h(U(F,m))\subseteq H(\mu)\},\\
 D_{k,m}&=\pi_2(Z_{k,m}),\qquad D_{0,m}=\varnothing.
 \end{split}
\]
The fixed-support estimate below is the analogue of
\cite[Claim~1, p.~651]{KrupskiMarciszewski2017}.  The linearly independent
set $P$ makes the present fibre finite.

\begin{proposition}\label{prop:closed-cover}
For all $k,m\in\N$, the set $Z_{k,m}$ is compact and $D_{k,m}$ is a
compact, hence closed, subset of $P$.  Moreover,
$D_{k-1,m}\subseteq D_{k,m}$ and
$P=\bigcup_{k,m\in\N}D_{k,m}$.
\end{proposition}

\begin{proof}
If $(F,\mu)\notin Z_{k,m}$, let us choose $f\in U(F,m)$ such that
$\abs{\mu(h(f))}>1$.
For every $x\in F$, continuity of $f$ gives an open neighbourhood
$V_x$ of $x$ on which $\abs f<1/m$.  Let
$V=\bigcup_{x\in F}V_x$.  Hence $f\in U(F',m)$ whenever
$F'\in\Fin_k(K)$ and $F'\subseteq V$.  The latter condition defines an
open Vietoris neighbourhood of $F$.  At the same time,
$\nu\mapsto\nu(h(f))$ is norm-continuous on $P$, so the strict inequality
persists on a relatively norm-open neighbourhood of $\mu$ in $P$.
These two neighbourhoods form a product
neighbourhood disjoint from $Z_{k,m}$.  Thus $Z_{k,m}$ is closed in the
compact Hausdorff space $\Fin_k(K)\times P$, and is compact.  Its second
projection $D_{k,m}$ is compact and hence closed in $P$.

The asserted inclusion follows from
$\Fin_{k-1}(K)\subseteq\Fin_k(K)$.  For the equality, let us fix $\mu\in P$.
Since $h(0)=0$, we have $\mu(h(0))=0$.  As $h$ is continuous into
$E_w$, the map
$f\mapsto\mu(h(f))$ is continuous at $0$ in $C_p(K)$.  There are a
finite set $F\subseteq K$
and $\varepsilon>0$ such that $\abs{f(x)}<\varepsilon$ for every
$x\in F$ implies $\abs{\mu(h(f))}<1$.
The set $F$ cannot be empty: otherwise surjectivity of $h$ would imply
$\abs{\mu(y)}<1$ for every $y\in E$, contrary to $\mu\ne0$.  Let us choose
$m\in\N$ with $1/m<\varepsilon$ and let $k=\abs F$.  Then
$(F,\mu)\in Z_{k,m}$, so $\mu\in D_{k,m}$.
\end{proof}


\begin{lemma}\label{lem:fixed-support}
For every $k,m\in\N$ and every $F\in\Fin_k(K)$ there are
$r\in\N$ and $\varphi_1,\ldots,\varphi_r\in E^*$, depending only on $(F,m)$,
such that
\[
 \{\mu\in P:(F,\mu)\in Z_{k,m}\}
 \subseteq P\cap\operatorname{span}\{\varphi_1,\ldots,\varphi_r\},
 \qquad
 \abs{\{\mu\in P:(F,\mu)\in Z_{k,m}\}}\leqslant r.
\]
\end{lemma}

\begin{proof}
Continuity of $h^{-1}$ at $0$ gives a weak neighbourhood of $0$ in $E$
of the form $W=\{y\in E:\abs{\varphi_i(y)}<\eta,
1\leqslant i\leqslant r\}$ such that
$h^{-1}(W)\subseteq U(F,m)$.  Since $h$ is bijective,
$W\subseteq h(U(F,m))$.

Let us fix $\mu\in P$ with $(F,\mu)\in Z_{k,m}$.  If
$\mu\notin\operatorname{span}\{\varphi_1,\ldots,\varphi_r\}$, there is
$y\in E$ such that
$\varphi_i(y)=0$ for $1\leqslant i\leqslant r$ and $\mu(y)=2$.
Indeed, for $T\colon E\to\R^r$, given by
$T(y)=(\varphi_1(y),\ldots,\varphi_r(y))$, if
$\ker T\subseteq\ker\mu$, the formula $\lambda(Ty)=\mu(y)$ defines a
linear functional on $T(E)$.  Since $T(E)$ is finite-dimensional,
$\lambda$ extends to a linear functional on $\R^r$.  Consequently,
$\mu=\lambda\circ T\in\operatorname{span}\{\varphi_1,\ldots,\varphi_r\}$.
Thus our assumption gives $y_0\in\ker T$ with $\mu(y_0)\ne0$; replacing
$y_0$ by $2y_0/\mu(y_0)$ gives the required vector $y$.

$y\in W\subseteq h(U(F,m))$, whereas
$h(U(F,m))\subseteq H(\mu)$.  This contradicts $\mu(y)=2$ and proves
the asserted inclusion.  The cardinal estimate follows from
Lemma~\ref{lem:arc}.
\end{proof}


\begin{thebibliography}{99}
\bibitem[AlbiacKalton2016]{AlbiacKalton2016} F.~Albiac and N.~J. Kalton, \emph{Topics in Banach Space Theory}, second ed., Graduate Texts in Mathematics, vol.~233, Springer, Cham, 2016.
\bibitem[Beezer2004]{Beezer2004} R.~A. Beezer, \emph{A First Course in Linear Algebra}, version~2.21, 2004, Section~VM, Theorem~DVM.
\bibitem[ErdosRado1960]{ErdosRado1960} P.~Erd\H{o}s and R.~Rado, \emph{Intersection theorems for systems of sets}, J. London Math. Soc. \textbf{35} (1960), no.~1, 85--90.
\bibitem[KrupskiMarciszewski2017]{KrupskiMarciszewski2017} M.~Krupski and W.~Marciszewski, \emph{On the weak and pointwise topologies in function spaces II}, J. Math. Anal. Appl. \textbf{452} (2017), 646--658.
\end{thebibliography}
\end{document}
